\chapter{Le fonti del diritto nazionale in materia di doping}

% ---------------------------------------------------------------------------
\section{Il quadro normativo nazionale}

Anche il \textbf{legislatore nazionale} si articola su due linee, che producono fonti distinte.

\begin{center}
\begin{forest} classalbero
[Legislatore nazionale
  [Stato
    [L.\ 376/2000]
  ]
  [NADO
    [NSA]
  ]
]
\end{forest}
\end{center}

% ---------------------------------------------------------------------------
\section{Le Norme Sportive Antidoping (NSA)}

\textbf{Norme Sportive Antidoping} (\textbf{NSA}): documento tecnico-attuativo del Codice
Mondiale Antidoping WADA e dei relativi Standard internazionali, adottato da NADO Italia. Si
articolano in tre atti:
\begin{itemize}
  \item \textbf{Codice Sportivo Antidoping} (\textbf{CSA}): attuativo del Codice WADA;
  \item \textbf{Disciplinare dei Controlli e delle Investigazioni} (\textbf{DC-I}): attuativo
    dell'\textit{International Standard for Testing and Investigations} WADA;
  \item \textbf{Disciplinare per le Esenzioni ai Fini Terapeutici} (\textbf{D-EFT}): in
    conformità con quanto previsto dall'\textit{International Standard for Therapeutic Use
    Exemptions}.
\end{itemize}

% ---------------------------------------------------------------------------
\section{L'ordinamento statale: la legge 376/2000}

\rubrica{Sanzioni penali}
\begin{itemize}
  \item \textbf{reclusione} (3 mesi -- 3 anni) e \textbf{multa} (5--100 milioni di lire) per il
    delitto di:
    \begin{itemize}
      \item procacciamento, somministrazione, assunzione o favoreggiamento dell'utilizzo di
        farmaci o sostanze dopanti;
      \item adozione o sottoposizione a pratiche mediche dopanti;
    \end{itemize}
  \item \textbf{reclusione} (2--6 anni) e \textbf{multa} (10--150 milioni) per il delitto di
    commercio illegale di farmaci e sostanze dopanti;
  \item \textbf{circostanze aggravanti} (danni, minore, dipendenti di organismi sportivi) e
    \textbf{pene accessorie} (interdizione dall'esercizio della professione).
\end{itemize}

% ---------------------------------------------------------------------------
\section{La definizione di doping tra i due ordinamenti}

\subsection{Ordinamento statale --- legge 376/2000}

\begin{quote}
\textit{``Costituiscono doping la somministrazione o l'assunzione di farmaci o di sostanze
biologicamente o farmacologicamente attive e l'adozione o la sottoposizione a pratiche mediche
non giustificate da condizioni patologiche ed idonee a modificare le condizioni psicofisiche o
biologiche dell'organismo al fine di alterare le prestazioni agonistiche degli atleti e comunque
idonee a modificare i risultati dei controlli sull'uso dei farmaci''}
\end{quote}

\subsection{Ordinamento sportivo --- il Codice WADA}

\begin{quote}
\textit{``Per doping si intende la violazione di una o più norme antidoping definite
dall'articolo 2.1 all'articolo 2.10 del Codice''}\\
\hfill --- Art.\ 1 del Codice WADA
\end{quote}

Le Norme Sportive Antidoping elencano dieci fattispecie di violazione (artt.\ 2.1--2.10):
\begin{enumerate}
  \item \textbf{presenza} di una sostanza vietata o dei suoi metaboliti o marker nel campione
    biologico dell'atleta (Art.\ 2.1);
  \item \textbf{uso o tentato uso} di una sostanza vietata o di un metodo proibito da parte di
    un atleta (Art.\ 2.2);
  \item \textbf{elusione o rifiuto} di sottoporsi al prelievo dei campioni biologici (Art.\ 2.3);
  \item \textbf{mancata reperibilità} (Art.\ 2.4);
  \item \textbf{manomissione o tentata manomissione} in relazione a qualsiasi fase dei controlli
    antidoping (Art.\ 2.5);
  \item \textbf{possesso} di sostanze vietate e metodi proibiti (Art.\ 2.6);
  \item \textbf{traffico} illegale o tentato traffico di sostanze vietate o metodi proibiti
    (Art.\ 2.7);
  \item \textbf{somministrazione o tentata somministrazione} ad un atleta durante le
    competizioni, di un qualsiasi metodo proibito o sostanza vietata, oppure somministrazione o
    tentata somministrazione ad un atleta, fuori competizione, di un metodo o di una sostanza che
    siano proibiti fuori competizione (Art.\ 2.8);
  \item \textbf{assistenza, incoraggiamento e aiuto, istigazione, complicità} in riferimento a
    una qualsiasi violazione o tentata violazione delle NSA (Art.\ 2.9);
  \item \textbf{divieto di associazione} (Art.\ 2.10).
\end{enumerate}

% ---------------------------------------------------------------------------
\section{L'Organizzazione Nazionale Antidoping (NADO)}

\textbf{NADO}: l'ente designato da ciascun Paese quale organismo investito dell'esclusiva
autorità e responsabilità in ordine a:
\begin{itemize}
  \item adozione e attuazione delle \textbf{norme antidoping};
  \item pianificazione e conduzione dei \textbf{controlli};
  \item \textbf{gestione dei risultati} delle analisi;
  \item svolgimento dei \textbf{processi}, in ambito nazionale;
  \item irrogazione delle \textbf{sanzioni sportive}.
\end{itemize}

Secondo la normativa WADA, in Italia il \textbf{C.O.N.I.} è la \textbf{N.A.D.O.} (Organizzazione
Nazionale Antidoping).

\rubrica{Requisiti richiesti dalla WADA alla NADO}
\begin{itemize}
  \item accettazione integrale del \textbf{Programma Mondiale Antidoping WADA}:
    \begin{itemize}
      \item Codice Mondiale Antidoping;
      \item Lista Sostanze e Metodi proibiti;
      \item Standard Internazionali (controlli, laboratori, esenzioni ai fini terapeutici);
      \item competenza del \textbf{TAS di Losanna};
    \end{itemize}
  \item certificazione di un sistema di qualità \textit{``UNI ISO 9001:2000''}.
\end{itemize}

% ---------------------------------------------------------------------------
\section{Strutture e funzioni di NADO Italia}

NADO Italia è organizzata in quattro strutture, ciascuna responsabile di una funzione:
\begin{itemize}
  \item \textbf{CCA}: pianificazione dei controlli;
  \item \textbf{PNA} (\textit{Procura Nazionale Antidoping}): gestione del risultato e delle
    indagini;
  \item \textbf{CEFT}: rilascio delle esenzioni ai fini terapeutici;
  \item \textbf{TNA}: organo giudicante.
\end{itemize}
